% !PRINTED 341
\weberpage{341}
%<<CONT
it follows that $a_{11}$ is divisible by $a_{22}$, and we set $a_{11} = a_{22}a$.
Then we have the basis of $\mathfrak{a}$:
\weq{8}{\begin{aligned} \alpha_1 &= a_{22}a,\\ \alpha_2 &= a_{12} + a_{22}\theta. \end{aligned}}
Then, by Bd.\,II, §\,164, (12):
\weq{9}{N(\mathfrak{a}) = a_{22}^2\,a,}
and if for $(1, \theta)$ we take the basis §\,90, (9), then
\[
\theta + \theta' = 0 \text{ or } = -1,
\]
\weq{10}{\theta\theta' = -\tfrac14\Delta \text{ or } = \frac{1-\Delta}{4}.}
According as $\Delta \equiv 0$, or $\equiv 1 \pmod 4$, and
\[
(\theta - \theta')^2 = \Delta.
\]
In these two cases there then results, when one substitutes (8) in (3)
and takes account of (9):
\[
a_{22}b = 2a_{12} \text{ or } = 2a_{12} - a_{22},
\]
and consequently in both cases:
\weq{11}{\alpha_1 = a_{22}a, \quad \alpha_2 = a_{22}\,\frac{b+\sqrt{\Delta}}{2},}
and the form
\[
T = (a,\,b,\,c)
\]
here has a \gesperrt{positive first coefficient}.

If we set
\weq{12}{\omega = -\frac{\alpha_2}{\alpha_1} = -\frac{b+\sqrt{\Delta}}{2a},}
then $\omega$ is a root of the quadratic equation:
\weq{13}{a\omega^2 + b\omega + c = 0}
and
\weq{14}{\lambda = a_{22}a(x-\omega y)}
is a functional corresponding to the ideal $\mathfrak{a}$, and at the same time a
basis form of $\mathfrak{a}$. $a$ is the least positive integer for which
$a\omega$ becomes an integer.

The necessary and sufficient condition for two
integers $\alpha_1, \alpha_2$ of the field $\Omega$ to form the basis of an ideal
consists in this, that first of all $\alpha_1, \alpha_2$ is a basis of the field $\Omega$,
and that, if $\omega_1, \omega_2$ is a minimal basis of this field,
\[
\alpha_1\omega_1,\quad \alpha_2\omega_1,\quad \alpha_1\omega_2,\quad \alpha_2\omega_2
\]
by the linear form
\[
\lambda = \alpha_1 t_1 + \alpha_2 t_2
\]
%CONT>>
